\documentclass[11pt,a4paper]{article}


\usepackage[utf8]{inputenc}
\usepackage[T1]{fontenc}
\usepackage[spanish,es-nodecimaldot]{babel}
\usepackage{lmodern}
\usepackage{microtype}
\usepackage{geometry}
\geometry{margin=2.25cm, headheight=22pt, headsep=0.55cm, footskip=1cm}
\usepackage{amsmath,amssymb}
\usepackage{booktabs,longtable,array,tabularx,colortbl}
\usepackage{enumitem}
\usepackage{xcolor}
\usepackage{tikz}
\usetikzlibrary{arrows.meta,positioning,shapes.geometric,fit,calc}
\usepackage{float}
\usepackage{pdflscape}
\usepackage{fancyhdr}
\usepackage{titlesec}
\usepackage{caption}
\usepackage{hyperref}

\definecolor{UHred}{HTML}{8F1D1D}
\definecolor{UHredDark}{HTML}{741515}
\definecolor{UHrose}{HTML}{FBF5F5}
\definecolor{UHline}{HTML}{B75A5A}
\hypersetup{colorlinks=true,linkcolor=UHred,urlcolor=UHred}
\titleformat{\section}{\Large\bfseries\color{UHredDark}}{\thesection.}{0.55em}{}
\titleformat{\subsection}{\large\bfseries\color{UHred}}{\thesubsection}{0.55em}{}
\titleformat{\subsubsection}{\normalsize\bfseries\color{UHred}}{\thesubsubsection}{0.55em}{}
\titlespacing*{\section}{0pt}{2.0ex plus .5ex}{1.0ex}
\captionsetup{font=small,labelfont={bf,color=UHredDark},justification=centering}
\setlength{\parskip}{3pt}
\setlength{\parindent}{1.1em}
\setlist{topsep=3pt,itemsep=2pt}
\renewcommand{\arraystretch}{1.12}

\pagestyle{fancy}
\fancyhf{}
\fancyhead[L]{\small\scshape Universidad de La Habana}
\fancyhead[R]{\small\scshape Arquitectura de Computadoras}
\fancyfoot[C]{\thepage}
\renewcommand{\headrulewidth}{0.5pt}
\renewcommand{\headrule}{\hbox to\headwidth{\color{UHline}\leaders\hrule height \headrulewidth\hfill}} \renewcommand{\footrulewidth}{0pt}

\title{\textbf{Diseño y funcionamiento de un microprocesador S-MIPS multiciclo}}
\author{Luis Enrique Cohalla Cabrera}
\date{16 de noviembre de 2026}

\arrayrulecolor{UHline}
\begin{document}
\begin{center}
    {\small\scshape\color{UHred} Universidad de La Habana $\cdot$ Arquitectura de Computadoras}\par
    \vspace{0.45em}
    {\LARGE\bfseries\color{UHredDark} Diseño y funcionamiento de un microprocesador S-MIPS multiciclo}\par
    \vspace{0.35em}
    {\normalsize Luis Enrique Cohalla Cabrera \quad -- \quad \textit{16 de noviembre de 2026}}\par
    \vspace{0.5em}
    {\color{UHline}\rule{\textwidth}{0.6pt}}
\end{center}
\vspace{0.5em}

\section{Introducción}
Este documento describe la arquitectura, el camino de datos y la unidad de control de un microprocesador compatible con el subconjunto S-MIPS implementado como una máquina multiciclo. La implementación utiliza una memoria compartida para instrucciones y datos, un banco de registros de 32 palabras de 32 bits, una ALU y una unidad de control basada en una máquina de estados finitos (FSM).
La ejecución de cada instrucción se divide en etapas. En cada ciclo, la FSM activa las señales de control necesarias para mover datos entre los registros internos, la ALU, la memoria y el banco de registros. El registro de instrucción conserva la instrucción durante su ejecución, mientras que el contador de programa (PC) se actualiza cuando corresponde.

\subsection{Convenciones}
\begin{itemize}[nosep]
    \item Todas las palabras y rutas principales de datos tienen 32 bits.
    \item El banco de registros contiene 32 registros de 32 bits.
    \item Los campos de instrucción se indican como \texttt{IR[31:0]}.
    \item Los inmediatos ordinarios se extraen de \texttt{IR[15:0]}.
    \item Las direcciones de memoria son direccionables por byte y la organización de palabras es little-endian. \item Los nombres de señales respetan mayúsculas y minúsculas tal como se usan en el circuito.
\end{itemize}

\section{Arquitectura general}

El procesador se organiza en los siguientes bloques:
\begin{enumerate}
    \item \textbf{PC}: contiene la dirección de la instrucción actual.
    \item \textbf{Memoria}: almacena instrucciones y datos. La dirección seleccionada procede del PC o del registro de dirección de memoria.
    \item \textbf{IR}: almacena la instrucción leída.
    \item \textbf{Register File}: banco de 32 registros con dos puertos de lectura y un puerto de escritura.
    \item \textbf{Registros A y B}: almacenan los operandos leídos del banco de registros durante la etapa de decodificación.
    \item \textbf{ALU}: realiza operaciones aritméticas, lógicas, comparaciones y desplazamientos.
    \item \textbf{ALUOut}: conserva el resultado de la ALU entre ciclos.
    \item \textbf{MDR}: conserva el dato leído de memoria.
    \item \textbf{Datapath}: contiene los multiplexores, extensores y registros que conectan los bloques anteriores.
    \item \textbf{Instruction Decoder, Control Unit y FSM}: decodifican la instrucción, determinan las señales de control y secuencian los ciclos.
    \item \textbf{Branch Control}: decide si una condición de salto condicional se cumple.
\end{enumerate}

\begin{figure}[H]
\centering
\begin{tikzpicture}[
        block/.style={draw=UHred,fill=UHrose,rounded corners=2pt,align=center,minimum width=2.3cm,minimum height=0.9cm,text=UHredDark},
        small/.style={draw=UHred,fill=UHrose,rounded corners=2pt,align=center,minimum width=1.8cm,minimum height=0.75cm,text=UHredDark},
        arr/.style={-{Latex[length=2mm]},thick,color=UHred},
        node distance=1.0cm and 1.35cm
]
\node[block] (pc) {PC};
\node[block,below=of pc] (mem) {Memoria};
\node[block,right=3.0 cm of pc] (ir) {IR};
\node[block,below=1.2cm of ir] (decoder) {Instruction\\Decoder};
\node[block,below=1.2cm of decoder] (control) {Control Unit};
\node[block,below=1.4cm of control] (rf) {Register File\\32 $\times$ 32};
\node[small,below=1.5cm of mem] (mdr) {MDR};
\node[small,below=1.5cm of mdr] (ab) {Registros\\A y B};
\node[block,below=1.4cm of ab] (alu) {ALU};
\node[small,left=1.3cm of alu] (aluout) {ALUOut};
\draw[arr] (pc) -- node[above]{dirección} (mem);
\draw[arr] (mem) -- node[above]{instrucción} (ir);
\draw[arr] (ir) -- (decoder);
\draw[arr] (decoder) -- (control);
\draw[arr] (control.west) -- ++(-1.0,0) |- (pc.east);
\draw[arr] (control.south) --  (rf.north);
\draw[arr] (rf) -- node[above]{RD1, RD2} (ab);
\draw[arr] (ab) -- (alu);
\draw[arr] (alu) -- (aluout);
\draw[arr] (mem.south) -- ++(0,-0.7) -| (mdr.north);
\draw[arr] (mdr) -- (rf);
\draw[arr] (aluout) |- (mem.west);
\draw[arr] (aluout.north) |- (pc.west);
\end{tikzpicture}
\caption{Vista funcional simplificada. Las flechas representan caminos de información; las señales de control se distribuyen desde la unidad de control.}
\end{figure}

\section{Formato de instrucción y decodificación}

El decodificador obtiene los campos de la instrucción:
\[
    op = IR[31:26],\quad rs=IR[25:21],\quad rt=IR[20:16],
\]
\[
    rd=IR[15:11],\quad shamt=IR[10:6],\quad funct=IR[5:0].
\]
Para instrucciones inmediatas se utiliza $imm=IR[15:0]$, y para saltos directos el campo de dirección procede de $IR[25:0]$.

El \textit{Instruction Decoder} genera señales de identificación como \texttt{IsAdd}, \texttt{IsSub}, \texttt{IsAnd}, \texttt{IsOr}, \texttt{IsSlt}, \texttt{IsSll}, \texttt{IsSrl}, \texttt{IsJr}, \texttt{IsSysCall}, \texttt{IsAddi}, \texttt{IsAndi}, \texttt{IsOri}, \texttt{IsLui}, \texttt{IsLw}, \texttt{IsSw}, \texttt{IsLb}, \texttt{IsLbu}, \texttt{IsSb}, \texttt{IsBeq}, \texttt{IsBne}, \texttt{IsJ}, \texttt{IsJal}, entre otras.

\subsection{Subconjunto de instrucciones considerado}

\begin{longtable}{@{}llp{8.2cm}@{}}
    \toprule
    \textbf{Clase} & \textbf{Instrucciones} & \textbf{Descripción} \\
    \midrule
    R & \texttt{add}, \texttt{addu}, \texttt{sub} & Suma y resta entre registros. \\
    R & \texttt{and}, \texttt{or} & Operaciones lógicas bit a bit. \\
    R & \texttt{slt} & Comparación con signo; produce 1 si $rs<rt$. \\
    R & \texttt{sll}, \texttt{srl} & Desplazamientos lógicos. \\
    R & \texttt{jr} & Salto a la dirección almacenada en $rs$. \\
    I & \texttt{addi}, \texttt{addiu} & Suma de registro e inmediato con extensión de signo. \\
    I & \texttt{slti} & Comparación con signo frente a inmediato extendido con signo. \\
    I & \texttt{andi}, \texttt{ori} & Operaciones lógicas con inmediato extendido con ceros. \\
    I & \texttt{lui} & Carga el inmediato en los 16 bits superiores. \\
    I & \texttt{lw}, \texttt{sw} & Lectura/escritura de palabra. \\
    I & \texttt{lb}, \texttt{lbu}, \texttt{sb} & Lectura de byte con extensión de signo o ceros, y escritura de byte.\\
    I & \texttt{beq}, \texttt{bne} & Saltos condicionales por igualdad o desigualdad. \\
    J & \texttt{j}, \texttt{jal} & Salto directo; \texttt{jal} guarda la dirección de retorno en $r31$. \\
    Sistema & \texttt{syscall} & Servicio de impresión de carácter o terminación. \\
    \bottomrule
\end{longtable}

\section{Camino de datos}

\subsection{Contador de programa y selección de dirección}

El PC se actualiza mediante el multiplexor \texttt{MUX\_PCSource}, de cuatro entradas y 32 bits:
\begin{center}
    \begin{tabularx}{0.9\textwidth}{cX}
        \toprule
        \textbf{PCSource} & \textbf{Dato seleccionado} \\
        \midrule
        00 & Registro A, para \texttt{jr}. \\
        01 & \texttt{ALUOut}, para destino de branch calculado. \\
        10 & \texttt{ALU\_RESULT}, normalmente $PC+4$. \\
        11 & \texttt{Jump\_Target}, para \texttt{j} y \texttt{jal}. \\
        \bottomrule
    \end{tabularx}
\end{center}

La dirección secuencial se calcula con la ALU:
\[
    PC+4 = PC + 4.
\]
El destino de salto condicional se calcula como:
\[
    PC_{\text{branch}} = (PC+4) + (\operatorname{SignExt}(imm)\ll 2).
\]

El destino de salto directo se construye como:
\[
    Jump\_Target = \{(PC+4)[31:28], IR[25:0], 2'b00\}.
\]

La escritura del PC está habilitada por:
\[
    PCEnable = PCwr \lor (PCwrCond \land BranchTaken).
\]
La unidad \texttt{Branch\_Control} calcula:
\[
    BranchTaken=(IsBeq\land ZeroFlag)\lor(IsBne\land\neg ZeroFlag).
\]

\subsection{Lectura del banco de registros}

El banco de registros tiene dos puertos de lectura y uno de escritura. Los campos $rs$ y $rt$ seleccionan los registros fuente. Las salidas se capturan en los registros internos A y B durante el estado S1.

El registro $r0$ (\texttt{\$zero}) siempre devuelve cero y no admite escritura efectiva. En el reinicio se inicializan:
\[
    \$gp=r28=0x00001800,\qquad \$sp=r29=0x00003FFC.
\]

La dirección de escritura y el dato de escritura se preparan mediante multiplexores:
\begin{itemize}
    \item \texttt{MUX\_RegisterDestination}: selecciona $rt$ o $rd$ según la clase de instrucción.
    \item \texttt{MUX\_JALDestination}: fuerza $r31$ cuando se ejecuta \texttt{jal}.
    \item \texttt{MUX\_WriteBack}: selecciona el dato de ALU o el dato leído de memoria.
    \item \texttt{MUX\_JALWriteBack}: selecciona $PC+4$ para escribir la dirección de retorno de \texttt{jal}.
\end{itemize}

\subsection{Selección de operandos de la ALU}
El multiplexor \texttt{MUX\_ALUSrcA} selecciona:
\begin{itemize}
    \item Entrada 0: PC.
    \item Entrada 1: registro A.
\end{itemize}

El multiplexor \texttt{MUX\_ALUSrcB} selecciona una de cuatro entradas:
\begin{center}
    \begin{tabularx}{0.9\textwidth}{cX}
        \toprule
        \textbf{ALUSrcB} & \textbf{Entrada seleccionada} \\
        \midrule
        00 & Constante 4. \\
        01 & Registro B. \\
        10 & Inmediato extendido (\texttt{ExtendedImm}). \\
        11 & Inmediato extendido con signo y desplazado dos bits. \\
        \bottomrule
    \end{tabularx}
\end{center}

\subsection{Extensión de inmediatos}

La implementación distingue tres transformaciones:
\begin{enumerate}
    \item \textbf{Extensión de signo}: copia IR[15] en los 16 bits superiores. Se utiliza en \texttt{addi}, \texttt{addiu}, \texttt{slti}, desplazamientos de branch y direcciones efectivas.
    \item \textbf{Extensión con ceros}: coloca IR[15:0] en los 16 bits inferiores y ceros en los superiores. Se utiliza en \texttt{andi} y \texttt{ori}.
    \item \textbf{LUI}: coloca el inmediato en los 16 bits superiores:
    \[
        LUI\_Value = \{IR[15:0],16'b0\}= \operatorname{ZeroExt}(imm)\ll16.
    \]
\end{enumerate}

Para elegir el inmediato ordinario se utiliza un multiplexor controlado por:
\[
    ZeroImmSelect = IsAndi \lor IsOri.
\]
Su entrada 0 es \texttt{SignExtendedImm}, su entrada 1 es \texttt{ZeroExtendedImm} y su salida se denomina \texttt{ExtendedImm}. Esta salida alimenta la entrada 2 de \texttt{MUX\_ALUSrcB}. La entrada de branch desplazada dos bits permanece separada.

\subsection{ALU y registro ALUOut}

La ALU recibe \texttt{ALUDataA} y \texttt{ALUDataB}. Su operación se codifica mediante \texttt{ALUControl[2:0]}:
\begin{center}
    \begin{tabular}{ccc l}
        \toprule
        \textbf{AC[2]} & \textbf{AC[1]} & \textbf{AC[0]} & \textbf{Operación} \\
        \midrule
        0 & 0 & 0 & ADD \\
        0 & 0 & 1 & SUB \\
        0 & 1 & 0 & AND \\
        0 & 1 & 1 & OR \\
        1 & 0 & 0 & SLT \\
        1 & 0 & 1 & SLL \\
        1 & 1 & 0 & SRL \\
        \bottomrule
    \end{tabular}
\end{center}

Las señales de selección de operación se forman a partir del estado y de las señales de identificación de instrucción:
\begin{align*}
    ALU_Add &= (IsS0\lor IsS1\lor IsS2)\\
        &\quad\lor(IsS6\land(IsAdd\lor IsAddu))\\
        &\quad\lor(IsS10\land(IsAddi\lor IsAddiu)),\\
    ALU_Sub &= IsS8\lor(IsS6\land IsSub),\\
    ALU_And &= (IsS6\land IsAnd)\lor(IsS10\land IsAndi),\\
    ALU_Or &= (IsS6\land IsOr)\lor(IsS10\land IsOri),\\
    ALU_Slt &= (IsS6\land IsSlt)\lor(IsS10\land IsSlti),\\
    ALU_Sll &= IsS6\land IsSll,\\
    ALU_Srl &= IsS6\land IsSrl.
\end{align*}

A partir de estas señales:
\[
AC[2]=ALU\_Slt\lor ALU\_Sll\lor ALU\_Srl,
\]
\[
AC[1]=ALU\_And\lor ALU\_Or\lor ALU\_Srl,
\]
\[
AC[0]=ALU\_Sub\lor ALU\_Or\lor ALU\_Sll.
\]

El resultado se captura en \texttt{ALUOut} cuando \texttt{ALUwr}=1. Para \texttt{lui}, el multiplexor \texttt{MUX\_LUI} selecciona el inmediato desplazado 16 bits en lugar del resultado ordinario de la ALU.

\subsection{Acceso a memoria y bytes}

La dirección de memoria se selecciona mediante \texttt{MUX\_MemoryAddress}:
\begin{itemize}
    \item PC durante la búsqueda de instrucción.
    \item \texttt{ALUOut} durante accesos a datos.
\end{itemize}
El registro MDR conserva el dato leído. Para las lecturas de byte, el circuito selecciona el byte correspondiente según los dos bits menos significativos de la dirección:
\begin{itemize}
    \item \texttt{lb}: extiende el bit de signo del byte leído.
    \item \texttt{lbu}: extiende con ceros.
\end{itemize}
Para \texttt{sb}, el dato de escritura se coloca en el carril de byte seleccionado y \texttt{MEM\_BE} habilita únicamente ese byte. En palabras, se habilitan los cuatro bytes. La memoria utiliza organización little-endian.

\section{Máquina de estados finitos}
La unidad de control es una FSM con un registro de estado de 4 bits. Se emplean los siguientes estados:

\begin{longtable}{@{}clp{9cm}@{}}
    \toprule
    \textbf{Código} & \textbf{Estado} & \textbf{Función} \\
    \midrule
    0000 & S0 -- FETCH & Selección del PC como dirección de memoria, lectura de instrucción y cálculo de $PC+4$. \\
    0001 & S1 -- DECODE & Decodificación, lectura del banco de registros y captura en A/B. \\
    0010 & S2 -- MEM\_ADDR & Cálculo de dirección efectiva para carga o almacenamiento. \\
    0011 & S3 -- MEM\_READ & Lectura de memoria y captura en MDR. \\
    0100 & S4 -- LOAD\_WB & Escritura del dato cargado en el banco de registros. \\
    0101 & S5 -- MEM\_WRITE & Escritura de palabra o byte en memoria. \\
    0110 & S6 -- R\_EXEC & Ejecución de operación tipo R. \\
    0111 & S7 -- R\_WB & Escritura del resultado tipo R. \\
    1000 & S8 -- BRANCH & Comparación y cálculo del destino de branch. \\
    1001 & S9 -- JUMP & Actualización del PC para salto, \texttt{jal} o \texttt{jr}. \\
    1010 & S10 -- IMM\_EXEC & Ejecución de operación aritmética/lógica inmediata o \texttt{lui}. \\
    1011 & S11 -- IMM\_WB & Escritura del resultado inmediato en el banco de registros. \\
    1100 & S12 -- SYSCALL & Decodificación del servicio de sistema solicitado. \\
    1101 & S13 -- PRINT & Emisión de un carácter por la salida TTY. \\
    1110 & S14 -- HALT & Detención permanente hasta reinicio. \\
    1111 & S15 -- RECOVERY & Estado de recuperación; retorna a FETCH. \\
    \bottomrule
\end{longtable}

\subsection{Diagrama de transiciones}

El diagrama muestra los dieciséis estados de la FSM y las condiciones que determinan cada transición. Las etiquetas sobre las flechas identifican la clase de instrucción, el servicio solicitado o el final de la operación.

\begin{landscape}
\thispagestyle{fancy}
\begin{figure}[H]
\centering
\begin{tikzpicture}[
    scale=0.98,transform shape,every node/.style={font=\small},
    st/.style={circle,draw=UHred,line width=1pt,fill=UHrose,text=UHredDark,minimum size=10mm,inner sep=1pt,font=\small\bfseries},
    dec/.style={diamond,draw=UHred,line width=0.9pt,fill=UHrose,text=UHredDark,align=center,aspect=2.0,inner sep=2pt,font=\scriptsize},
    done/.style={rounded corners=2pt,draw=UHred,fill=UHrose,text=UHredDark,align=center,font=\scriptsize,inner sep=3pt},
    arr/.style={-{Latex[length=2.2mm]},draw=UHred,line width=0.9pt},
    lab/.style={fill=white,inner sep=2pt,text=black,font=\scriptsize}
]
% Estados de la FSM.
\node[st] (s0) at (-7,0) {$S_0$};
\node[st] (s1) at (-4,0) {$S_1$};
\node[dec] (decode) at (-1,0) {Clase de\\instrucción};
\node[st] (s2) at (2,2) {$S_2$};
\node[st] (s3) at (5,2) {$S_3$};
\node[st] (s4) at (8,2) {$S_4$};
\node[st] (s5) at (8,0.2) {$S_5$};
\node[st] (s6) at (2,0) {$S_6$};
\node[st] (s7) at (5,0) {$S_7$};
\node[st] (s8) at (3,-2.8) {$S_8$};
\node[st] (s9) at (5,-1.8) {$S_9$};
\node[st] (s10) at (1.5,-3.6) {$S_{10}$};
\node[st] (s11) at (5,-5.3) {$S_{11}$};
\node[st] (s12) at (-1,3.5) {$S_{12}$};
\node[dec] (service) at (2,3.5) {Servicio};
\node[st] (s13) at (6,3.2) {$S_{13}$};
\node[st] (s14) at (6,4.4) {$S_{14}$};
\node[st] (s15) at (-4 , -2.3) {$S_{15}$};
\node[done] (finish) at (10,-2.2) {Fin de\\operación};

% Inicio y decodificación.
\node[left=8mm of s0,font=\small] {Start / Reset};
\draw[arr] (-8.1,0) -- (s0);
\draw[arr] (s0) -- node[lab,above]{Always} (s1);
\draw[arr] (s1) -- node[lab,above]{Decode} (decode);

% Despacho según la instrucción.
\draw[arr] (decode) to[out=35,in=215] node[lab,above left]{IsLoad $\lor$ IsStore} (s2.west);
\draw[arr] (decode) -- node[lab,above]{IsR} (s6.west);
\draw[arr] (decode) to[out=-55,in=160] node[lab,below left]{IsBeq $\lor$ IsBne} (s8.west);
\draw[arr] (decode) to[out=-35,in=145] node[lab,below]{IsJ / IsJr} (s9.west);
\draw[arr] (decode) to[out=-70,in=170] node[lab,below]{IsImmediate} (s10.west);
\draw[arr] (decode) to[out=120,in=180] node[lab,left]{Syscall} (s12.west);
\draw[arr] (decode.south west) to[out=300,in=90] node[lab,left]{Illegal} (s15.north);


% Rutas de memoria, tipo R e inmediatas.
\draw[arr] (s2) -- node[lab,above]{IsLoad} (s3);
\draw[arr] (s2) to[bend right=5] node[lab,right]{IsStore} (s5);
\draw[arr] (s3) -- node[lab,above]{Always} (s4);
\draw[arr] (s6) -- node[lab,above]{Always} (s7);
\draw[arr] (s10) to[bend right=30] node[lab,above]{Always} (s11);

% Syscall y selección del servicio.
\draw[arr] (s12) -- (service);
\draw[arr] (service) to[bend right=12] node[lab,above]{\$v0 = 11} (s13);
\draw[arr] (service) to[bend left=12] node[lab,right]{\$v0 = 10} (s14);
\draw[arr] (s14) edge[loop right] node[right,font=\scriptsize]{HALT / permanece} (s14);
\draw[arr] (s12.north west) to[out=160,in=90] node[lab,left]{Servicio no reconocido} (s0.north);

% Todas las rutas que terminan regresan a FETCH por el punto común Done.
\draw[arr] (s4.east) to[out=0,in=120] node[lab,above]{Done} (finish.north);
\draw[arr] (s5.east) to[out=0,in=180] node[lab,above]{Done} (finish.west);
\draw[arr] (s7.east) to[out=0,in=200] node[lab,below]{Done} (finish.west);
\draw[arr] (s8.south) to[out=-90,in=200] node[lab,below]{Done} (finish.south west);
\draw[arr] (s9.south) to[out=-90,in=220] node[lab,below]{Done} (finish.west);
\draw[arr] (s11.east) to[out=0,in=270] node[lab,below]{Done} (finish.south west);
\draw[arr] (s13.east) to[bend left=60,out=-300,in=90] node[lab,right]{Done / TTY\_WE} (finish.north);
\draw[arr] (finish.south) to[out=-90,in=-90] node[lab,below]{Next Fetch} (s0.south);

% Recuperación desde instrucción ilegal.
\draw[arr] (s15) -- node[lab,left]{Recovery} (s0);
\end{tikzpicture}
\caption{Diagrama completo de estados y transiciones de la FSM S-MIPS (codificación de 4 bits).}
\end{figure}
\end{landscape}

\subsection{Condiciones de transición}

\begin{longtable}{@{}clp{8.5cm}@{}}
    \toprule
    \textbf{Estado actual} & \textbf{Siguiente estado} & \textbf{Condición} \\
    \midrule
    S0 & S1 & Siempre. \\
    S1 & S2 & Instrucción de carga o almacenamiento. \\
    S1 & S6 & Instrucción tipo R, excepto saltos y syscall. \\
    S1 & S8 & \texttt{beq} o \texttt{bne}. \\
    S1 & S9 & \texttt{j}, \texttt{jal} o \texttt{jr}. \\
    S1 & S10 & Instrucción inmediata, incluida \texttt{lui}. \\
    S1 & S12 & \texttt{syscall}. \\
    S1 & S15 & Instrucción ilegal o no reconocida, según la lógica de decodificación. \\
    S2 & S3 & Carga (\texttt{lw}, \texttt{lb}, \texttt{lbu}). \\
    S2 & S5 & Almacenamiento (\texttt{sw}, \texttt{sb}). \\
    S3 & S4 & Siempre. \\
    S4, S5, S7, S8, S9, S11, S13 & S0 & Fin de la operación correspondiente. \\
    S6 & S7 & Siempre. \\
    S10 & S11 & Siempre. \\
    S12 & S13 & Servicio de impresión: \texttt{\$v0}=11. \\
    S12 & S14 & Servicio de terminación: \texttt{\$v0}=10. \\
    S12 & S0 & Servicio no reconocido. \\
    S14 & S14 & Permanece detenido hasta reset. \\
    S15 & S0 & Recuperación al estado inicial de búsqueda. \\
    \bottomrule
\end{longtable}

\subsection{Diagrama temporal de una instrucción}

Como ejemplo, una instrucción de carga de palabra \texttt{lw} recorre:
\[
    S0 \rightarrow S1 \rightarrow S2 \rightarrow S3 \rightarrow S4 \rightarrow S0.
\]
Una instrucción tipo R recorre:
\[
S0 \rightarrow S1 \rightarrow S6 \rightarrow S7 \rightarrow S0.
\]
Una instrucción inmediata recorre:
\[
    S0 \rightarrow S1 \rightarrow S10 \rightarrow S11 \rightarrow S0.
\]
Una instrucción de almacenamiento recorre:
\[
    S0 \rightarrow S1 \rightarrow S2 \rightarrow S5 \rightarrow S0.
\]

\section{Señales de control}

Las señales se generan combinando la decodificación de estado con las señales de instrucción.

\subsection{Señales principales por estado}

En la tabla, una expresión como $S0\lor S9$ indica que la señal se activa en cualquiera de esos estados. Las señales no listadas en una fila permanecen desactivadas salvo que se indique lo contrario.

\begin{longtable}{@{}lp{10cm}@{}}
    \toprule
    \textbf{Señal} & \textbf{Ecuación de control} \\
    \midrule
    \texttt{PCwr} & $S0 \lor S9$ \\
    \texttt{PCwrCond} & $S8$ \\
    \texttt{IRwr} & $S0$ \\
    \texttt{IorD} & $S3 \lor S5$ \\
    \texttt{MDRwr} & $S3$ \\
    \texttt{Memwrite} & $S5$ \\
    \texttt{Regwr} & $S4 \lor S7 \lor S11 \lor (S9\land IsJal)$ \\
    \texttt{ALUwr} & $S1 \lor S2 \lor S6 \lor S10$ \\
    \texttt{Awr}, \texttt{Bwr} & $S1$ \\
    \texttt{TTY\_WE} & $S13$ \\
    \texttt{HALT} & $S14$ \\
    \texttt{RegDst} & $S4 \lor S11$ \\
    \texttt{MemToReg} & $S4$ \\
    \texttt{ALUSrcA} & $S2 \lor S6 \lor S8 \lor S10$ \\
    \texttt{ALUSrcB[1]} & $S1 \lor S2 \lor S10$ \\
    \texttt{ALUSrcB[0]} & $S1 \lor S6 \lor S8$ \\
    \texttt{PCSource[1]} & $S0 \lor S1 \lor (S9\land\neg IsJr)$ \\
    \texttt{PCSource[0]} & $S8 \lor (IsS9\land\neg IsJr)$ \\
    \bottomrule
\end{longtable}

\subsection{Interpretación de PCSource}

\begin{center}
    \begin{tabular}{ccl}
        \toprule
        \textbf{Estado / caso} & \textbf{PCSource} & \textbf{Resultado} \\
        \midrule
        S0 & 10 & $PC+4$ \\
        S8 & 01 & Destino de branch en ALUOut \\
        S9 con \texttt{jr} & 00 & Registro A ($rs$) \\
        S9 con \texttt{j}/\texttt{jal} & 11 & Jump\_Target \\
        \bottomrule
    \end{tabular}
\end{center}

\section{Instrucciones de salto y enlace}

\begin{itemize}
    \item \texttt{beq}: se calcula $A-B$. Si \texttt{ZeroFlag}=1, se habilita la escritura del PC con el destino de branch.
    \item \texttt{bne}: se calcula $A-B$. Si \texttt{ZeroFlag}=0, se habilita la escritura del PC con el destino de branch.
    \item \texttt{j}: el PC recibe \texttt{Jump\_Target}.
    \item \texttt{jal}: el PC recibe \texttt{Jump\_Target} y el banco de registros recibe $PC+4$ en $r31$.
    \item \texttt{jr}: el PC recibe el contenido de $rs$, almacenado en A.
\end{itemize}

\section{Llamadas al sistema}

La instrucción \texttt{syscall} se atiende en S12. El circuito selecciona explícitamente los registros de argumentos:
\begin{itemize}
    \item El puerto de lectura correspondiente al servicio selecciona $r2$ (\texttt{\$v0}).
    \item El puerto de lectura del argumento selecciona $r4$ (\texttt{\$a0}).
\end{itemize}

Se comparan los valores:
\[
    \$v0=11 \Rightarrow \text{servicio de impresión},
    \qquad
    \$v0=10 \Rightarrow \text{terminación}.
\]
Para el servicio 11, S12 pasa a S13. En S13 se activa \texttt{TTY\_WE} y se entrega a la salida TTY el carácter formado por los 7 bits inferiores de \texttt{\$a0}. Para el servicio 10, S12 pasa a S14 y la señal \texttt{HALT} permanece activa mientras la FSM se mantiene en S14.

\section{Secuencias de control resumidas}

\begin{table}[H]
    \centering
    \begin{tabularx}{\textwidth}{lX}
        \toprule
        \textbf{Instrucción} & \textbf{Secuencia de estados} \\
        \midrule
        \texttt{add}, \texttt{sub}, \texttt{and}, \texttt{or}, \texttt{slt}, \texttt{sll}, \texttt{srl}, \texttt{addu} & S0 $\to$ S1 $\to$ S6 $\to$ S7 $\to$ S0 \\
        \texttt{addi}, \texttt{addiu}, \texttt{slti}, \texttt{andi}, \texttt{ori}, \texttt{lui} & S0 $\to$ S1 $\to$ S10 $\to$ S11 $\to$ S0 \\
        \texttt{lw}, \texttt{lb}, \texttt{lbu} & S0 $\to$ S1 $\to$ S2 $\to$ S3 $\to$ S4 $\to$ S0 \\
        \texttt{sw}, \texttt{sb} & S0 $\to$ S1 $\to$ S2 $\to$ S5 $\to$ S0 \\
        \texttt{beq}, \texttt{bne} & S0 $\to$ S1 $\to$ S8 $\to$ S0 \\
        \texttt{j}, \texttt{jal}, \texttt{jr} & S0 $\to$ S1 $\to$ S9 $\to$ S0 \\
        \texttt{syscall} de impresión & S0 $\to$ S1 $\to$ S12 $\to$ S13 $\to$ S0 \\
        \texttt{syscall} de salida & S0 $\to$ S1 $\to$ S12 $\to$ S14 $\to$ S14 $\to\cdots$ \\
        \bottomrule
    \end{tabularx}
    \caption{Recorridos de la FSM para las principales clases de instrucciones.}
\end{table}

\section{Inicialización y reinicio}

Al activar el reset, el procesador vuelve al estado inicial de búsqueda y se inicializan los registros especiales indicados en el diseño:
\[
    r28 = 0x00001800,\qquad r29 = 0x00003FFC.
\]
El registro $r0$ conserva siempre el valor cero. El estado S14 no sale por transición normal: requiere reinicio para volver a ejecutar instrucciones.

\section{Verificación mediante programas de prueba}

La validación se realiza ejecutando programas de ensamblador y comparando la salida TTY con la salida de referencia de MARS.

\section{Conclusión}

El procesador implementa una arquitectura multiciclo en la que una FSM de 4 bits coordina el camino de datos. La separación entre búsqueda, decodificación, ejecución, acceso a memoria y escritura permite reutilizar la ALU y la memoria en distintos ciclos. Los multiplexores controlan las fuentes de operandos, el destino del PC y la selección del dato que se escribe en el banco de registros. El decodificador identifica las instrucciones y la unidad de control traduce esa identificación en señales temporales.

\end{document}